\defmodule {PRIM}

This module provides some procedures to search for integers
$m$ that are prime, for which $r = (m^k-1)/(m-1)$ is also prime
for certain values of $k$, and perhaps for which $(m-1)/2$ is also
prime.  The values of $r$ are always represented as {\tt SuperInteger}'s.

\bigskip\hrule
\code

DEFINITION MODULE PRIM;

FROM SUPINT   IMPORT SuperInteger;
FROM SUPFACT  IMPORT Factors;


PROCEDURE Trier (VAR L : Factors);
\endcode
 \tab 
  Sorts the factors in the list {\tt L} in {\em decreasing\/} order.
 \endtab
\code

PROCEDURE SFindPrimer (k, e : LONGINT;  S1, S2 : SuperInteger;  
                       Safe : BOOLEAN);
\endcode
 \tab
  Finds and prints all integers $m$ such that $S_1 \le m \le S_2$ and    
  such that both $m$ and $r = (m^k-1)/(m-1)$ are prime.   
  If {\tt Safe = TRUE}, it is also required that $(m-1) / 2$ be prime.
  If $k=1$, then $r$ is considered prime.  
  The retained $m$ are printed in decreasing order,
  in the form $m = 2^e \pm x$ and also as decimal integers. 
  The factors of $m-1$ are also printed.
 \endtab
\code

PROCEDURE FindPrimer (k, m0, m1 : LONGINT;  Safe : BOOLEAN);
\endcode
 \tab
  Similar to {\tt SFindPrimemr}, but for values of $m$ less than 
  $2^{31}$.
 \endtab
\code

PROCEDURE FindPrimesmr 
      (k, e, c1, c2 : LONGINT;  Safe : BOOLEAN;   F : ARRAY OF CHAR);
\endcode
 \tab
   Calls {\tt SFindPrimer} for $S_1 = 2^e + c_1$ and $S_2 = 2^e + c_2$, 
   creates the file {\tt F} and prints the results in it.
 \endtab
\code

END PRIM.
\endcode


